% Section III: research gap and objectives. Included in main.tex after the literature review.
\section{Research Gap and Objectives}
\label{sec:gap}
Four things are missing from the work above.
\emph{G1:} as far as we can tell, nobody has compared the tool definitions inside package archives across
versions. Kraishan looked only at registry
fields~\cite{kraishan2026drift}, so the rate at which honest servers change tools, and therefore the price of a tool pin,
is unknown.
\emph{G2:} every existing tool check reads the tools from a running server~\cite{invariant2025mcpscan,mcppin,mcpseal}.
By the time it complains, the new code has already run.
\emph{G3:} no check compares a release with the one before it. Scanners look at one version~\cite{sejfia2022amalfi},
provenance names a builder without remembering the previous one~\cite{slsa2023provenance}, and lockfiles keep a
checksum~\cite{gamage2025lockfiles}. A stolen account that publishes from a new workflow, under a new user or with a new
install script gets past all of them.
\emph{G4:} the pinning tools were never run against real attacks, and the pre-install hook was tested on 11 hand-made
scenarios~\cite{bagmar2026setup}. A detection rate means little without the number of honest upgrades that a person must
then unblock by hand.

The four objectives map onto these gaps. O1 measures tool drift in MCP packages and checks whether metadata
changes with it (G1, RQ1); without that number nobody can say what a tool pin costs. O2 is the artifact itself, a
lockfile plus a pre-launch gate that never runs package code (G2, G3). O3 tests the gate against the closest earlier
one on real attacks and benign upgrades, under a kill rule registered in advance~\cite{osf2026prereg} (G3, G4, RQ2).
In O4 an attacker gets full knowledge of the gate, and each fix is priced on benign upgrades (G2, G3, RQ3).
Fig.~\ref{fig:method} shows the method.
